\begin{frame}{발표: 5$\sim$7분의 설계도}
  \begin{center}
    \small
    \begin{tabular}{@{}lcl@{}}
      \toprule
      순서 & 내용 & 시간 \\
      \midrule
      1 & 문제 정의 & 1분 \\
      2 & 데이터와 전처리 & 1분 \\
      3 & 적용한 모델과 이유 & 2분 \\
      4 & 결과와 수식적 정당화 & 2분 \\
      5 & 잘 안 됐던 부분과 원인 진단 & 1분 \\
      \bottomrule
    \end{tabular}
  \end{center}
  \vskip0.5em
  \begin{itemize}
    \item ``정확도 90\% 달성''보다 \textbf{``왜 이 모델을 골랐고
      왜 이 결과가 나왔는지''}를 설명하는 쪽에 높은 점수.
    \item 마지막 1분(``잘 안 됐던 부분'')이 \textbf{가장
      평가받는} 부분 --- rubric의 ``결과의 정직성과 반성''이
      정확히 여기 를 본다.
    \item ``모든 것이 잘 됐다''는 발표는 이 항목에서 \textbf{스스로
      감점}. ML2의 Week 16.1(강화학습 캡스톤)도 동일한 구조.
  \end{itemize}
\end{frame}
